%% Finitary Jónsson--Tarski representation and descriptive duality.

\section{Boolean algebras with operators}

\begin{annexitem}{annex:p10.bao.1}{thm:p10-bao-ultrafilter-witness}{Simultaneous ultrafilter witnesses}
\begin{proof}
Binary-join preservation implies coordinatewise monotonicity:
if \(a\leq b\), then \(a\vee b=b\), and additivity gives
the value at \(a\) below the value at \(b\).
Fix \(U\) and the filters in the statement, with \(n>0\).
We enlarge one coordinate filter at a time to an ultrafilter,
retaining the condition that every tuple from the current filters
has its \(f\)-value in \(U\).

At stage \(k\), the current filters are
\(V_1,\ldots,V_{k-1},F_k,\ldots,F_n\).
Let \(I_k\) consist of those \(b\in B\) for which some choices
in every other current coordinate give an \(f\)-value outside \(U\)
when \(b\) is placed in coordinate \(k\).
We show that \(I_k\) is an ideal.
It contains \(0\): choose \(1\) in every other coordinate and use
normality. It is downward closed by monotonicity.

Suppose \(b,d\in I_k\), with respective witness tuples
\((c_i)_{i\ne k}\) and \((e_i)_{i\ne k}\).
Put \(t_i=c_i\wedge e_i\) in every other coordinate, including
coordinates whose filters have not yet been enlarged.
All \(t_i\) remain in their current filters.
Monotonicity puts both values \(f(\ldots,b,\ldots)\) and
\(f(\ldots,d,\ldots)\), evaluated at these common \(t_i\),
below their respective failed witness values.
Neither belongs to \(U\), since an element of an upward-closed
filter cannot lie below an element outside it.
Additivity in coordinate \(k\) and ultrafilter join-primeness give
\[
 f(\ldots,b\vee d,\ldots)
   =f(\ldots,b,\ldots)\vee f(\ldots,d,\ldots)\notin U.
\]
Thus the same tuple witnesses \(b\vee d\in I_k\).

The induction invariant implies \(I_k\cap F_k=\varnothing\).
The filter generated by \(F_k\) and the complements of \(I_k\) is
\[
 E_k=\{z:\exists a\in F_k\ \exists i\in I_k\
                                a\wedge\neg i\leq z\}.
\]
It contains the generators and \(1\), is upward closed, and is
closed under meets: combine two filter witnesses by meet and two
ideal witnesses by join.
It is proper, because \(a\wedge\neg i=0\) would imply \(a\leq i\)
and hence \(a\in I_k\), contradicting disjointness.
BPIT extends \(E_k\) to an ultrafilter \(V_k\).
This ultrafilter contains \(F_k\) and excludes \(I_k\).
Consequently every tuple from
\(V_1,\ldots,V_k,F_{k+1},\ldots,F_n\) still has value in \(U\).
Finite induction proves the filter form of the lemma.

For the element form, \(f(a_1,\ldots,a_n)\in U\) and normality
imply that every \(a_i\ne0\).
Use the proper principal filters \(\{b:a_i\leq b\}\).
Monotonicity verifies their tuple condition, so the filter form
gives the required witnesses.
For \(n=0\), no enlargement is needed; the conclusion is exactly
the premise. When \(B\) is trivial, there is no \(U\), so the
assertion is vacuous.
Every instance uses finitely many BPIT extensions.
The relations for a set of operator symbols are defined by their
universal clauses, with no simultaneous selection of witnesses.
\end{proof}
\end{annexitem}

\begin{annexitem}{annex:p10.bao.2}{thm:p10-jonsson-tarski}{J\'onsson--Tarski representation}
\begin{proof}
First, each complex operation is normal and preserves arbitrary
unions in every positive-arity coordinate.
Membership of a successor coordinate in a union means membership
in at least one member of that union, exactly the existential
condition in the complex-operation definition.
For the empty union there is no witness, giving normality.
The nullary operation is its designated unary relation extension.

Stone representation makes \(\beta_B\) an injective Boolean
homomorphism, with image all clopens of \(X_B\).
For \(n>0\), if \(U\) belongs to
\(f^{(\mathbf B^\bullet)^+}(\widehat b_1,\ldots,\widehat b_n)\),
choose the witnessing \(V_i\).
They contain \(b_i\), so their defining relation implies
\(f^B(b_1,\ldots,b_n)\in U\).
This proves inclusion in \(\widehat{f^B(b_1,\ldots,b_n)}\).
Conversely, membership in this latter set satisfies the element
form of the ultrafilter-witness lemma.
Its \(V_i\) lie in the required basic clopens and satisfy \(R_f\),
proving membership in the complex operation.
For \(n=0\), the equality follows immediately from
\(R_f(U)\) iff \(f^B()\in U\).
Every operation is therefore preserved.
The Boolean image is closed under those operations by the same
equalities.
The trivial algebra maps isomorphically to the power set of the
empty frame, with its uniquely possible operations.
\end{proof}
\end{annexitem}

\begin{annexitem}{annex:p10.bao.3}{thm:p10-bao-canonical-completion}{The dense and compact canonical completion}
\begin{proof}
\paragraph{The equivalent compactness formulations.}
Suppose the finite-witness condition holds for subsets of \(e[B]\).
Let \((k_i)_{i\in I}\) be closed elements and
\((o_j)_{j\in J}\) open elements, with
\(\bigwedge_i k_i\leq\bigvee_j o_j\).
A closed element is the meet of all embedded Boolean elements
above it, and an open element is the join of all embedded Boolean
elements below it: any one defining meet or join verifies the
reverse inequality to the immediate bound.
Put
\[
 S=\bigcup_{i\in I}\{b:k_i\leq e(b)\},\qquad
 T=\bigcup_{j\in J}\{b:e(b)\leq o_j\}.
\]
Then \(\bigwedge_i k_i=\bigwedge e[S]\) and
\(\bigvee_j o_j=\bigvee e[T]\).
Compactness gives finite \(S_0,T_0\) witnessing the inequality.
For each of these finitely many Boolean elements select one
index witnessing its membership in the corresponding union.
The finite sets \(I_0,J_0\) of selected indices satisfy
\[
 \bigwedge_{i\in I_0} k_i
   \leq\bigwedge e[S_0]
   \leq\bigvee e[T_0]
   \leq\bigvee_{j\in J_0}o_j.
\]
Empty selected sets use the bounds \(1\) and \(0\).
Conversely, every embedded element is both closed and open, so
the family formulation implies the image-subset formulation.
This establishes equivalence without any infinite choice of
representations.

\paragraph{Density and compactness of the Stone embedding.}
The power-set algebra is complete and atomic, with singleton atoms.
For every \(U\in X_B\),
\[
 \{U\}=\bigcap_{a\in U}\widehat a,\qquad
 X_B\setminus\{U\}=\bigcup_{a\notin U}\widehat a.
\]
For the first equality, an ultrafilter containing all of \(U\)
equals \(U\) by maximality.
For the second, distinct ultrafilters have an element belonging
to the other ultrafilter and absent from \(U\).
Every subset is a union of its singletons and an intersection
of the complements of the singletons outside it.
These two representations prove density.

Suppose \(\bigcap_{s\in S}\widehat s
\subseteq\bigcup_{t\in T}\widehat t\).
If no finite \(S_0,T_0\) satisfies
\(\bigwedge S_0\leq\bigvee T_0\), then every finite meet from
\(S\cup\{\neg t:t\in T\}\) is nonzero.
The filter it generates is proper and extends by BPIT to an
ultrafilter \(U\).
This \(U\) contains all of \(S\) and all complements of \(T\),
contradicting the displayed inclusion.
The finite inequality therefore exists.
Its image is the required finite inequality in the completion.
Empty meets and joins are included throughout.

\paragraph{Atoms in an arbitrary dense compact completion.}
Let \(e:B\to D\) be another such Boolean completion.
For \(U\in X_B\), put
\[
 p_U=\bigwedge_{a\in U}e(a).
\]
Compactness implies \(p_U\ne0\).
Indeed, \(p_U\leq e(0)\) would yield a finite subset of \(U\)
whose meet has zero image; injectivity of \(e\) would then put
zero in the proper filter \(U\).

For every \(a\), ultrafilter dichotomy gives
\(p_U\leq e(a)\) or \(p_U\leq\neg e(a)\).
If a closed element \(k=\bigwedge e[S]\) is nonzero and
\(k\leq p_U\), every \(s\in S\) must belong to \(U\);
otherwise \(k\leq e(s)\) and \(k\leq p_U\leq\neg e(s)\)
would force \(k=0\).
Thus \(p_U\leq k\), so \(k=p_U\).
Join-density by closed elements now shows \(p_U\) is an atom:
every nonzero element below it contains a nonzero closed element,
which must be \(p_U\).

Every nonzero \(d\in D\) contains a nonzero closed element
\(k=\bigwedge e[S]\), again by density.
The set
\[
 F_k=\{a\in B:k\leq e(a)\}
\]
is a proper filter. It contains \(S\), and
\(\bigwedge e[F_k]=k\).
Extend it to an ultrafilter \(U\) by BPIT.
Then \(0<p_U\leq k\leq d\).
Thus \(D\) is atomic.
Moreover, its atoms are precisely the \(p_U\):
any atom contains one of them and therefore equals it.
Distinct ultrafilters give distinct atoms, since an element
separating the ultrafilters puts the two atoms below complementary
Boolean elements.

\paragraph{The unique isomorphism over \(B\).}
Define
\[
 j:D\longrightarrow\mathcal P(X_B),\qquad
 j(d)=\{U:p_U\leq d\}.
\]
The atom laws show that this preserves finite Boolean operations
and both bounds.
It is injective because a nonzero Boolean difference contains
an atom. It is onto: for \(Y\subseteq X_B\), take
\(d=\bigvee_{U\in Y}p_U\).
An atom outside this family lies below the complement of each
member, so their join lies below its complement; hence it does
not lie below \(d\). This proves \(j(d)=Y\).

For each \(a\in B\), \(p_U\leq e(a)\) holds exactly when
\(a\in U\), so \(j(e(a))=\widehat a\).
Any isomorphism over \(B\) is an order isomorphism and preserves
all existing meets and joins.
It must therefore send \(p_U=\bigwedge e[U]\) to
\(\bigcap_{a\in U}\widehat a=\{U\}\).
Every element of a complete atomic Boolean algebra is the join
of the atoms below it: a nonzero difference from that join would
contain a further atom.
The isomorphism is consequently determined everywhere by its
values on the \(p_U\), proving uniqueness.
For trivial \(B\), preservation of \(0,1\) forces \(D\) to be
trivial too, and the assertion holds directly.

The statements about the operators of \(\mathbf B^\sigma\) follow
from their relational definitions and the arbitrary-union
calculation in Annex~\ref{annex:p10.bao.2}.
No uniqueness of an arbitrary operator extension is inferred merely
from uniqueness of the underlying Boolean completion.
\end{proof}
\end{annexitem}

\begin{annexitem}{annex:p10.bao.4}{thm:p10-bao-descriptive-representation}{Descriptive representation and reconstruction}
\begin{proof}
\paragraph{The general frame of an algebra.}
For \(\mathbf B_*\), the admissible sets are the Stone clopens
and are closed under the complex operators by representation.
They separate points.
A family of these clopens with nonempty finite intersections
has nonempty total intersection by compactness of the Stone
space, using the open-complement cover argument.
The frame is tight because its defining canonical relation says
exactly that every tuple of admissible neighborhoods has its
operator value containing the root point.
The nullary instance is the same defining membership clause.
Thus the general frame is descriptive, and the representation
isomorphism onto its admissible algebra follows from
Theorem~\ref{thm:p10-jonsson-tarski} and Stone surjectivity.

\paragraph{The topology of a descriptive frame.}
Conversely, let \(\mathcal G=(W,(R_f),\mathcal A)\) be descriptive.
The admissible sets form a clopen basis, since they form a Boolean
algebra containing \(W\).
Differentiation gives the Hausdorff property: an admissible set
separating two points and its complement are disjoint neighborhoods.
For an arbitrary open cover, take all admissible sets subordinate
to one of its members.
If the cover has no finite subcover, the complements of this
whole refinement have nonempty finite intersections.
Otherwise finitely many refinement members would cover \(W\),
and finitely many parent cover members would do so.
Descriptive compactness now gives a point in all the complements,
contradicting the refinement being a cover.
Thus the topology is compact.

Every clopen is a finite union of members of this basis, by
compactness applied to its basis cover together with its open
complement. Finite unions belong to \(\mathcal A\).
Consequently \(\mathcal A=\operatorname{Clop}(W)\).
The empty space satisfies these conclusions as well: its
admissible algebra has one member.

\paragraph{The reconstruction map preserves the relations.}
Stone reconstruction makes
\[
 \varepsilon(w)=\{A\in\mathcal A:w\in A\}
\]
a homeomorphism \(W\to X_{\mathcal A}\) which identifies the
admissible families.
For each \(f\), the canonical relation on the right holds at
\(\varepsilon(w),\varepsilon(w_1),\ldots,\varepsilon(w_n)\)
exactly when every choice of admissible \(A_i\) containing \(w_i\)
has \(w\in f^{\mathcal G^+}(A_1,\ldots,A_n)\).
By tightness this is equivalent to the original
\(R_f(w,w_1,\ldots,w_n)\).
For arity zero the same equivalence contains no chosen neighborhoods.
Thus \(\varepsilon\) preserves and reflects all relations,
completing the general-frame isomorphism.
\end{proof}
\end{annexitem}

\begin{annexitem}{annex:p10.bao.5}{thm:p10-bao-descriptive-duality}{Duality with descriptive general frames}
\begin{proof}
\paragraph{A general-frame morphism induces a homomorphism.}
Let \(g:\mathcal G\to\mathcal H\) be such a morphism.
Its inverse-image map is defined on the admissible algebras and
preserves their Boolean operations.
For admissible \(A_1,\ldots,A_n\) of \(\mathcal H\), the back
condition gives one direction, and the forth condition the other,
of
\[
 g^{-1}[f^{\mathcal H^+}(A_1,\ldots,A_n)]
   =f^{\mathcal G^+}(g^{-1}[A_1],\ldots,g^{-1}[A_n]).
\]
Indeed, a successor tuple at \(g(w)\) lifts by back to a tuple
at \(w\); a tuple at \(w\) maps by forth to one at \(g(w)\).
Membership in the coordinate sets is preserved by these
preimage conditions.
For a nullary symbol the equality is its required relation
equivalence. Thus the map is a homomorphism.

\paragraph{The dual of a homomorphism has the forth property.}
Let \(h:\mathbf A\to\mathbf B\) be a homomorphism and
\(g(U)=h^{-1}[U]\).
Stone duality proves that this maps ultrafilters to ultrafilters
and that preimages of basic clopens are basic clopens:
\[
 g^{-1}[\widehat a]=\widehat{h(a)}.
\]
If \(R_f^{\mathbf B^\bullet}(U,V_1,\ldots,V_n)\) and
\(a_i\in g(V_i)\), then \(h(a_i)\in V_i\), so
\[
 h(f^A(a_1,\ldots,a_n))
       =f^B(h(a_1),\ldots,h(a_n))\in U.
\]
This is the required canonical relation at the image tuple.

\paragraph{The back property for all positive arities.}
Suppose \(R_f^{\mathbf A^\bullet}(g(U),W_1,\ldots,W_n)\),
with \(n>0\).
For each coordinate define the image-generated filter
\[
 F_i=\{b\in B:\exists a\in W_i\quad h(a)\leq b\}.
\]
It contains \(1\), is upward closed, and is closed under meets.
It is proper in this situation.
For if \(a\in W_i\) and \(h(a)=0\), use \(a\) in that coordinate
and \(1\) in all others in the assumed source relation.
It gives \(f^A(1,\ldots,a,\ldots,1)\in g(U)\).
Applying \(h\), operator preservation and normality give \(0\in U\),
which is impossible.
Thus no generator is zero and \(0\notin F_i\).

For any tuple \(b_i\in F_i\), select finitely many witnesses
\(a_i\in W_i\) with \(h(a_i)\leq b_i\).
The assumed relation and preservation by \(h\) give
\(f^B(h(a_1),\ldots,h(a_n))\in U\).
Monotonicity then gives \(f^B(b_1,\ldots,b_n)\in U\).
The filter form of
Lemma~\ref{thm:p10-bao-ultrafilter-witness} therefore supplies
ultrafilters \(V_i\supseteq F_i\) with
\(R_f^{\mathbf B^\bullet}(U,V_1,\ldots,V_n)\).
Each \(g(V_i)\) is an ultrafilter containing \(W_i\), since
\(h[W_i]\subseteq F_i\subseteq V_i\).
Maximality gives \(g(V_i)=W_i\), proving back.
For nullary \(f\), both properties are exactly
\[
 f^A()\in h^{-1}[U]
       \quad\Longleftrightarrow\quad
 h(f^A())=f^B()\in U.
\]
No successor choices are needed.

\paragraph{Functor laws and reconstruction.}
Inverse images preserve identities and reverse composition.
The defined morphisms are closed under composition: forth composes
directly, back lifts a tuple successively through the two maps,
and admissible preimages compose.
Thus both constructions are contravariant functors.
The preceding descriptive-representation theorem gives
isomorphisms on objects whose underlying maps are \(\beta\)
and \(\varepsilon\) from Stone duality.
They respect operators, relations, and admissible sets by that theorem.

Their naturality and the two inverse identities are the Stone
membership calculations of Annex~\ref{annex:p10.stone.4},
now applied to the admissible algebras.
Explicitly, for a homomorphism \(h:\mathbf A\to\mathbf B\),
the inverse-image homomorphism induced by \(h_*\) sends
\(\widehat a\) to \(\widehat{h(a)}\).
For \(g:\mathcal G\to\mathcal H\), the inverse image of the
point ultrafilter at \(w\) under the induced algebra map is
the point ultrafilter at \(g(w)\).
These are the naturality equations, and the same calculations
with the reconstruction maps give identity composites.
Hence the two functors give a contravariant equivalence.
Empty objects and nullary relations obey the same identities.
\end{proof}
\end{annexitem}
